% !TEX TS-program = xelatex
% !TEX encoding = UTF-8 Unicode
% !Mode:: "TeX:UTF-8"

\documentclass{resume}
\usepackage{zh_CN-Adobefonts_external} 
\usepackage{linespacing_fix} % disable extra space before next section
\usepackage{cite}
\usepackage{fontawesome}
\usepackage{multicol}
\usepackage{hyperref}
\usepackage{xcolor} % Add xcolor for color support
\definecolor{crimson}{RGB}{200,0,0} % Define crimson color

\hypersetup{
hidelinks,
colorlinks=true,
allcolors=black,
pdfstartview=Fit,
breaklinks=true
}

\begin{document}
\pagenumbering{gobble}

\name{\textbf{吴世涵}}  % 强化姓名
\vspace{5pt}

\contactInfo{电子科技大学}{计算机科学与工程学院}{计算机科学与技术}{}
\vspace{5pt}

\centerline{\sffamily\large{
\faHome\ \href{https://koorye.github.io}{koorye.github.io}
\textperiodcentered\faLink\ \href{https://koorye.github.io/blog}{koorye.github.io/blog}
\textperiodcentered\faEnvelope\ \href{mailto:shihan.wu.koorye@outlook.com}{shihan.wu.koorye@outlook.com}
\textperiodcentered\faPhone\ +8615869663967
}}

\vspace{10pt}
\section{\faUser\ 个人介绍}

我是来自电子科技大学\textbf{计算机科学与技术}专业的学术硕士(学院排名\textbf{Top 1.3\%})，研究方向聚焦：

\faAngleRight \textbf{视觉语言模型}的高效与高性能迁移学习 \\
\faAngleRight \textbf{机器人视觉-语言-动作模型(VLA)}的训练与测试时适应

以\textbf{一作/共一}身份发表\textbf{CVPR(CCF-A)}论文2篇，另有2篇在投，获得\textbf{国家奖学金}、\textbf{优秀毕业生}等荣誉

\section{\faGraduationCap\ 教育背景}
\datedline{\textbf{电子科技大学}\quad 计算机科学与技术 · 学术硕士}{\textit{2023.9 $\sim$ 2026.6}}
\faAngleRight\textbf{学业表现}：\textit{GPA 3.94/4.0} \quad|\quad \textit{学院排名：6/454 \textcolor{crimson}{(Top 1.3\%)}} \\
\faAngleRight\textbf{核心荣誉}：\textbf{\textcolor{crimson}{国家奖学金}}、\textbf{优秀研究生}、学业一等奖学金、新生一等奖学金等

\vspace{5pt}
\datedline{\textbf{电子科技大学}\quad 软件工程 · 工学学士}{\textit{2019.9 $\sim$ 2023.6}}
\faAngleRight\textbf{学业表现}：\textit{GPA 3.94/4.0} \quad|\quad \textit{专业排名：18/181 (Top 10\%)} \quad|\quad \textit{英语四六级：579/467} \\
\faAngleRight\textbf{荣誉奖项}：\textbf{\textcolor{crimson}{优秀毕业生}}、“世强”专项奖学金、优秀学生奖学金等

\section{\faLightbulbO\ 科研论文}
\datedline{
   \textbf{[ICLR 2025 投稿中·共同第一作者]}\
   \textit{Policy Contrastive Decoding for Robotic Foundation Models}
}
{\textit{2025.9}}
\faAngleRight 针对\textbf{具身智能大模型}(\textit{VLA})的\textbf{虚假相关性}问题，首次提出\textbf{策略对比解码}的解决方法

\faAngleRight 通过\textbf{开放语义目标检测模型}\textit{Grounding DINO}和\textbf{图像分割模型}\textit{SAM2}对图像观察中的目标实现实时分割和跟踪，并通过\textbf{图像重绘框架}\textit{Inpaint Anything}抹除目标物体得到负样本

\faAngleRight 通过正负样本在\textit{VLA}输出概率空间中的\textbf{对比解码}技术实现对虚假相关性的有效抑制

\faAngleRight 适用于基于\textbf{自回归}和\textbf{扩散}的多种\textit{VLA}策略，\textbf{无需训练}实现多种主流\textit{VLA}\textbf{+8\%$\sim$41\%}的性能提升

\faAngleRight 项目主页：\href{https://koorye.github.io/proj/PCD}{\textcolor{blue}{\sffamily https://koorye.github.io/proj/PCD}}，ArXiv：\href{https://arxiv.org/abs/2505.13255}{\textcolor{blue}{\sffamily https://arxiv.org/abs/2505.13255}}

\vspace{2pt}
\datedline{
   \textbf{[ICRA 2025 投稿中·共同第一作者]}\
   \textit{InSpire: Vision-Language-Action Models with Intrinsic Spatial Reasoning}
}
{\textit{2025.9}}
\faAngleRight 针对\textbf{具身智能大模型}(\textit{VLA})的\textbf{虚假相关性}问题，提出基于\textbf{空间推理}和\textbf{思维链}技术的解决方法

\faAngleRight 引入空间推理问题作为多模态大模型(\textit{VLM})与具身智能大模型(\textit{VLA})的中间桥梁，引导模型的动作推理关注图像观察中的目标

\faAngleRight 实现\textit{VLA}在可见\textbf{(+6.2\%)}和未见\textbf{(+10\%)}任务性能显著提升

\faAngleRight 项目主页：\href{https://koorye.github.io/proj/Inspire}{\textcolor{blue}{\sffamily https://koorye.github.io/proj/Inspire}}，ArXiv：\href{https://arxiv.org/abs/2505.13888}{\textcolor{blue}{\sffamily https://arxiv.org/abs/2505.13888}}

\vspace{5pt}
\datedline{
   \textbf{[\textcolor{crimson}{CVPR 2025 (CCF-A)}·第一作者]}\
   \textit{Skip Tuning: Pre-trained Vision-Language Models are Effective and Efficient Adapters Themselves}
}{\textit{2024.12}}
\faAngleRight 传统提示调优(\textit{P-Tuning}等)和适配器(\textit{LoRA}等)需要引入额外参数，且时间和内存效率提升有限

\faAngleRight 提出了针对\textbf{视觉语言模型}全新的\textbf{高效微调}方法，从长度和宽度两个维度修剪特征-梯度传播流

\faAngleRight 不改变结构的同时，实现\textbf{+1.04\%准确率}以及\textbf{15x时间效率}和\textbf{6.4x内存效率}的提升

\faAngleRight 项目主页：\href{https://github.com/Koorye/SkipTuning}{\textcolor{blue}{\sffamily https://github.com/Koorye/SkipTuning}}，ArXiv：\href{https://arxiv.org/abs/2412.11509}{\textcolor{blue}{\sffamily https://arxiv.org/abs/2412.11509}}

\vspace{5pt}
\datedline{
   \textbf{[\textcolor{crimson}{CVPR 2024 (CCF-A)}·共同第一作者]}\
   \textit{DePT: Decoupled prompt tuning}}
{\textit{2023.11}}
\faAngleRight 针对现有\textbf{提示调优}方法在已见和未见类上的性能权衡问题，提出\textbf{视觉语言模型}的\textbf{解耦学习}方法

\faAngleRight 将视觉语言模型的输出特征空间分为任务特定和任务通用空间，在学习已见类知识的同时保持未见类上的通用泛化能力

\faAngleRight 作为\textbf{即插即用}的解耦方法，取得现有多种提示调优方法\textbf{+0.67\%$\sim$2.65\%的性能普遍提升}

\faAngleRight 项目主页：\href{https://github.com/Koorye/DePT}{\textcolor{blue}{\sffamily https://github.com/Koorye/DePT}}，ArXiv：\href{https://arxiv.org/abs/2309.07439}{\textcolor{blue}{\sffamily https://arxiv.org/abs/2309.07439}}

\vspace{-2pt}
\section{\faCoffee\ 实习经历}

\datedline{\textbf{北京智源人工智能研究院（BAAI）}\quad 研究实习生·具身智能大模型}{\textit{2025.6 - 2025.10}}
\faAngleRight 设计\textbf{大型双臂交互式机器人数据集}，包含\textbf{6}款双臂机器人本体和\textbf{350K+}轨迹数据

\faAngleRight 探索\textbf{具身智能大模型}(VLA)的\textbf{能力金字塔}，设计双臂复杂任务的\textbf{层次化技能标注}

\faAngleRight 开发\textbf{Realman RM75, Agilex Piper, Pika}等双臂机器人的部署系统，以及基于\textbf{LeRobot}的数据平台

\section{\faBriefcase\ 项目经历}
\datedline{\textbf{雷达信号智能检测系统} · 算法研发}{\textit{2024.3 – 2024.7}}
\faAngleRight 设计雷达信号的实时检测系统，利用频域图实时检测并提取其中的雷达信号片段

\faAngleRight 改进一维残差网络作为信号\textbf{特征提取网络}和\textbf{分类器}，通过\textbf{数据增强}以实现更好的分类性能

\faAngleRight 设计\textbf{开集检测算法}，实现对已知和未知雷达信号的有效区分

\vspace{5pt}
\datedline{\textbf{面向自动驾驶场景的复杂目标检测系统} · 算法研发}{\textit{2023.1 – 2023.6}}
\faAngleRight 改进两阶段\textbf{目标检测模型}\textit{Faster R-CNN}实现未知物体的小样本适应

\faAngleRight 针对华为\textit{CODA}数据集中的极端场景，利用开放语义\textbf{视觉语言模型}\textit{CLIP}生成语义密集的软标签

\faAngleRight 利用\textbf{知识蒸馏}技术对\textit{Faster R-CNN}的\textit{RPN Head}与分类器提供引导，实现未知物体的有效检测

\vspace{5pt}
\datedline{\textbf{水面目标智能检测系统} · 算法研发}{\textit{2022.7 – 2022.9}}
\faAngleRight 选用单阶段\textbf{目标检测模型}\textit{YOLOX-l}，实现水面目标的实时检测

\faAngleRight 针对长尾数据分布设计了\textit{Mosaic}与\textit{Copy-Paste}\textbf{数据增强}策略，提升困难目标的检测性能

\vspace{5pt}
\datedline{\textbf{智能安全驾驶系统} · 后端开发与运维}{\textit{2022.4 – 2022.8}}
\faAngleRight 面向安全驾驶监测需求，设计了一套\textbf{嵌入式设备、服务端与检测算法}结合的智能安全驾驶系统

\faAngleRight 通过摄像头和可穿戴式设备实时收集用户图像和心跳、血压等信息，与服务端实时同步

\faAngleRight 调用检测算法对用户的疲劳驾驶、吸烟等危险行为进行监控和反馈

\vspace{5pt}
\datedline{\textbf{“FACE ME”人脸美学评分系统与智能美妆推荐系统} · 算法研发}{\textit{2021.10 – 2022.4}}
\faAngleRight 基于华南理工\textit{SCUT-FBP5500}人脸美学数据集，利用\textbf{残差网络}\textit{ResNet-18}回归预测对人脸进行评分

\faAngleRight 设计基于个性化\textit{PageRank}算法智能产品\textbf{推荐系统}，分析用户购买记录网络，为用户推荐感兴趣产品

% 专利与竞赛经历分类强化
\section{\faCogs\ 专利成果}

\datedline{\textbf{[申请中·学生第二发明人]}\quad 视觉语言模型的微调方法}{\textit{2025.2}}
\datedline{\textbf{[申请中·学生第一发明人]}\quad 基于低频增强的小样本图像分类迁移方法}{\textit{2024.8}}

\section{\faTrophy\ 竞赛奖项}

\datedline{\textbf{国际级Meritorious Winner}·美国大学生数学建模大赛(MCM)}{\textit{2021.4}}
\datedline{\textbf{国家级铜奖}·中国国际“互联网+”大学生创新创业大赛}{\textit{2021.12}}
\datedline{\textbf{省部级一等奖}·中国成都国际软件设计和应用竞赛}{\textit{2021.10}}
\datedline{\textbf{省部级二等奖}·全国大学生数学建模竞赛}{\textit{2021.9}}
\datedline{\textbf{省部级银奖}·“挑战杯”四川省大学生创业计划竞赛}{\textit{2022.7}}

% 技术栈分层展示
\section{\faCode\ 技术能力}

\textbf{核心能力}: 视觉语言模型(CLIP、LLaVa)及微调(Prompt Tuning、LoRA)、VLA模型(OpenVLA、$\pi_0$等)

\textbf{编程语言}: Python、Java、C、C\#、JavaScript、SQL等

\textbf{算法框架}: PyTorch、Keras、MMDetection、Scikit-learn等

\textbf{工程能力}: 前后端(Vue、SpringBoot)、数据库(MySQL、Redis)、运维(Docker)、游戏(Unity3D)、移动端

\end{document}